\title{\textbf{Progressive Token Pruning: Learnable Token Selection for Efficient Transformer Attention}

\author{Anonymous authors}



\section{Introduction}

\subsection{Motivation}

Transformer inference at long sequence length is dominated by the 
$\mathcal{O}(L^2)$ cost of softmax attention, but the actual \textbf{useful} 
content in a sequence is often much smaller than $L$. In classification 
tasks like SST-2, only a handful of tokens -- the sentiment-bearing words -- 
determine the output. In question answering, only tokens near the answer 
span are critical. The rest are redundant.

Prior work exploits this in three ways:

\begin{enumerate}
\item \textbf{Fixed heuristics.} Prune tokens by attention magnitude (Kim \& Cho,
  2021), or by position (longformer-style). Simple but not task-aligned.

\item \textbf{Learned selectors.} Train a selector network on top of a frozen
  transformer (Wang et al., 2021). Requires task-specific training.

\item \textbf{Full retraining.} Train the transformer with pruning in the loop
  (Rao et al., 2021). Effective but expensive.

\end{enumerate}
\textbf{The gap.} None of these is a \textbf{drop-in} solution that (a) learns which 
tokens to keep, (b) preserves pretrained weights exactly, and (c) works 
across flat, nested, and causal architectures. PTP fills this gap.

\subsection{Problem Statement}

Given a pretrained transformer with hidden dimension $d$ and $H$ heads, we 
seek a modified attention module $\mathcal{A}'$ that:

\begin{enumerate}
\item \textbf{Prunes tokens:} Retains only a fraction $\rho = k/L$ of tokens.
\item \textbf{Preserves [CLS]:} The classification token is always retained.
\item \textbf{Handles padding:} Invalid tokens are excluded.
\item \textbf{Preserves causality:} For decoders, no future information leaks.
\item \textbf{Drop-in:} Pretrained projections are copied exactly.
\item \textbf{Quality:} $\Delta\text{Acc} \leq 0.02\%$ on SST-2.

\end{enumerate}
\subsection{Our Approach: PTP}

\textbf{Progressive Token Pruning} (PTP) introduces a **learnable scoring 
network**:

\[
s_i = \sigma\left(\text{MLP}(x_i)\right) \in (0, 1)
\tag{1.1}
\]

where $x_i$ is the hidden state of token $i$, and the MLP is:

\[
\text{MLP}(x) = W_2 \cdot \text{GELU}(W_1 x + b_1) + b_2
\tag{1.2}
\]

with $W_1 \in \mathbb{R}^{64 \times d}$, $W_2 \in \mathbb{R}^{1 \times 64}$.

\textbf{Top-$k$ selection.} Let $k = \lceil \rho L \rceil$. We keep the $k$ 
tokens with the highest scores:

\[
\mathcal{T} = \text{top-}k_i \, s_i
\tag{1.3}
\]

\textbf{Key design choices:}

\begin{enumerate}
\item \textbf{[CLS] preservation.} We set $s_0 = 10^4$ before top-$k$, guaranteeing
  the first token is always retained.
\item \textbf{Padding masking.} Invalid tokens get $s_i = -10^4$ before top-$k$.
\item \textbf{Gather-based KV pruning.} Rather than masking, we gather the top-$k$
  KV tensors, reducing memory.
\item \textbf{Causal-aware attention.} For decoders, we mask queries that would
  attend to future keys by comparing original positions.

\end{enumerate}
\subsection{Contributions}

\begin{itemize}
\item \textbf{A learnable token-selection module.} PTP scores tokens with a small
  MLP and prunes them via top-$k$.
\item \textbf{[CLS] preservation.} Guaranteed by score clamping, with no additional
  parameters or logic.
\item \textbf{Causal-aware pruning.} We preserve causality by comparing original
  positions after pruning.
\item \textbf{A drop-in wrapper.} Pretrained projections are copied exactly; only
  the score network and LoRA adapters are trainable.
\item \textbf{A systematic evaluation.} Benchmarked on DistilBERT, RoBERTa, and
  GPT-2 on SST-2 with ablations over keep ratio and score network depth.
\item \textbf{Token importance analysis.} We show that learned scores correlate with
  content-word importance and reveal systematic patterns.

\end{itemize}
\subsection{Paper Organization}

Section 2 reviews related work. Section 3 provides background on attention 
and token pruning. Section 4 presents PTP. Section 5 describes integration 
with four architectures. Section 6 details the experimental setup. Section 
7 reports results. Section 8 analyzes the learned token scores. Section 9 
discusses practical guidance. Section 10 lists limitations. Section 11 
concludes.



\subsection{Notation}

\begin{table}[H]
\centering
\small
\begin{tabular}{ll}
\toprule
\textbf{Symbol} & \textbf{Meaning} \\
\midrule
$L$ & Sequence length \\
$H$ & Number of heads \\
$d$ & Hidden dimension \\
$\rho$ & Keep ratio \\
$k$ & Number of retained tokens ($= \lceil \rho L \rceil$) \\
$s_i$ & Score of token $i$ \\
$\mathcal{T}$ & Set of retained token indices \\
$x_i$ & Hidden state of token $i$ \\
$\text{MLP}$ & Score network \\
\bottomrule
\end{tabular}
\end{table}




\section{Related Work}

\subsection{Token Pruning}

\textbf{PoWER-BERT} (Goyal et al., 2020): Progressive word elimination based on 
attention scores. Fixed heuristic.

\textbf{TR-BERT} (Ye et al., 2021): Reinforcement learning for token selection. 
Requires RL training.

\textbf{LTP} (Kim \& Cho, 2021): Learned threshold pruning. Task-specific.

\textbf{SparseBERT} (Xia et al., 2022): Distilled sparsity patterns.

\textbf{PTP differs} by using a \textbf{learnable MLP} for scoring, retaining 
\textbf{[CLS] unconditionally}, and being \textbf{drop-in} for pretrained models.

\subsection{Attention Sparsity}

\textbf{Longformer} (Beltagy et al., 2020): Fixed window + global tokens.

\textbf{BigBird} (Zaheer et al., 2020): Adds random connections.

\textbf{Sparse Transformer} (Child et al., 2019): Strided patterns.

\textbf{PTP differs} by pruning \textbf{tokens} (dense per remaining token), not 
\textbf{connections}.

\subsection{Token Merging}

\textbf{Token Merging} (Bolya et al., 2022): Merges similar tokens via bipartite 
matching.

\textbf{PTP differs} by \textbf{hard pruning} rather than merging.

\subsection{Head Pruning}

\textbf{Voita et al. (2019):} Prune attention heads. Redundant heads dropped.

\textbf{Michel et al. (2019):} Single head suffices for many tasks.

\textbf{PTP operates orthogonally} -- it prunes tokens, not heads. Can be combined.

\subsection{Position of Our Work}

\begin{table}[H]
\centering
\small
\begin{tabular}{lllll}
\toprule
\textbf{Axis} & \textbf{PoWER-BERT} & \textbf{TR-BERT} & \textbf{LTP} & \textbf{PTP} \\
\midrule
Learnable &  & (RL) &  & (supervised) \\
 {[}CLS] preserved &  &  &  &  \\
Drop-in &  &  &  &  \\
Multi-architecture &  &  &  &  \\
Causal-aware &  &  &  &  \\
\bottomrule
\end{tabular}
\end{table}

\subsection{Comparison Table}

\begin{table}[H]
\centering
\small
\begin{tabular}{lllll}
\toprule
\textbf{Method} & \textbf{Selector} & \textbf{Loss} & \textbf{Causal} & \textbf{Pretrained} \\
\midrule
PoWER-BERT & Attention & CE &  &  \\
TR-BERT & RL policy & RL &  &  \\
LTP & Threshold & CE &  &  \\
\textbf{PTP} & \textbf{MLP} & \textbf{CE} & \textbf{ } & \textbf{ } \\
\bottomrule
\end{tabular}
\end{table}




\section{Background}

\subsection{Dense Softmax Attention}

For $Q, K, V \in \mathbb{R}^{L \times d_h}$:

\[
\text{Attn}(Q, K, V) = \text{softmax}\!\left(\frac{Q K^\top}{\sqrt{d_h}}\right) V \in \mathbb{R}^{L \times d_h}
\tag{3.1}
\]

\textbf{Complexity.} $\mathcal{O}(L^2 d_h)$ time, $\mathcal{O}(L^2)$ memory.

\subsection{Token Pruning}

Restrict attention to a subset $\mathcal{T} \subseteq \{1, \dots, L\}$ 
of \textbf{retained tokens}:

\[
\text{Attn}_\text{pruned}(Q, K, V) = \text{softmax}\!\left(\frac{Q_\mathcal{T} K_\mathcal{T}^\top}{\sqrt{d_h}}\right) V_\mathcal{T}
\tag{3.2}
\]

where subscript $\mathcal{T}$ denotes indexing.

\textbf{Complexity.} $\mathcal{O}(L |\mathcal{T}| d_h)$ time, $\mathcal{O}(L |\mathcal{T}|)$ memory.

For $|\mathcal{T}| = \rho L$:

\[
\text{Complexity reduction} = \rho
\tag{3.3}
\]

\subsection{Scoring Network}

A learnable function $f_\theta: \mathbb{R}^d \to \mathbb{R}$ maps hidden 
states to scores. Choices:

\begin{table}[H]
\centering
\small
\begin{tabular}{lll}
\toprule
\textbf{Type} & \textbf{Expression} & \textbf{Params} \\
\midrule
Linear & $w^\top x$ & $d$ \\
MLP-1 & $W_2 \cdot \text{GELU}(W_1 x)$ & $64 d$ \\
MLP-2 & Two hidden layers & More \\
\bottomrule
\end{tabular}
\end{table}

PTP uses \textbf{MLP-1} (one hidden layer of 64 units). This is a good 
trade-off between expressiveness and parameter count.

\subsection{Top-K Selection}

Given scores $s \in \mathbb{R}^L$:

\[
\mathcal{T} = \text{top-}k(s)
\tag{3.4}
\]

\textbf{Complexity.} $\mathcal{O}(L \log k)$ via heap or $\mathcal{O}(L)$ via 
quickselect.

\subsection{[CLS] Preservation}

In classification tasks, the [CLS] token (position 0) carries the 
aggregated sequence representation. Dropping it breaks classification.

\textbf{PTP guarantees [CLS] retention} by setting $s_0 = 10^4$ before top-$k$:

\[
s_0 \leftarrow 10^4
\tag{3.5}
\]

This is a soft guarantee (no runtime branching required).

\subsection{Padding Handling}

Padding tokens have zero or invalid content. They must be excluded:

\[
s_i \leftarrow -10^4 \quad \text{if token } i \text{ is padding}
\tag{3.6}
\]

\subsection{Causal Masking}

For decoder models, position $i$ cannot attend to $j > i$. After pruning, 
we must preserve this by comparing \textbf{original} positions:

\[
\text{mask}_{ij} = \mathbb{1}[\text{orig}(j) > i]
\tag{3.7}
\]

\subsection{Rotary Position Embeddings}

RoPE (Su et al., 2021) applies position-dependent rotation:

\[
\tilde{Q}_i = Q_i \cos(\theta_i) + \text{rot}(Q_i) \sin(\theta_i)
\tag{3.8}
\]

PTP applies RoPE to Q, K before pruning.

\subsection{Notation Summary}

\begin{table}[H]
\centering
\small
\begin{tabular}{ll}
\toprule
\textbf{Symbol} & \textbf{Meaning} \\
\midrule
$L$ & Sequence length \\
$\rho$ & Keep ratio \\
$k$ & Retained tokens \\
$s_i$ & Score of token $i$ \\
$\mathcal{T}$ & Retained indices \\
$f_\theta$ & Score network \\
\bottomrule
\end{tabular}
\end{table}




\section{Method}

\subsection{Overview}

PTP is a wrapper for standard attention that replaces the full $L \times L$ 
attention with a pruned $L \times k$ attention:

\begin{enumerate}
\item \textbf{Score} each token with a small MLP.
\item \textbf{Select} the top-$k$ tokens.
\item \textbf{Prune} K, V to the selected tokens.
\item \textbf{Attend} using the pruned K, V.

\end{enumerate}
\textbf{Figure 4.1} illustrates the pipeline.

\begin{lstlisting}
    X $\in$ R^{B $$\times$ L $$\times$ d}
          |
          +------------------+
                              
   +--------------+    +-------------+
   | Score MLP    |    | Q/K/V proj  |
   | s = $\sigma$(MLP(X))|    +------+------+
   +------+-------+           |
                              |
   +--------------+           |
   | Top-k        |           |
   | T = top_k(s) |           |
   +------+-------+           |
          |                   |
                               
   +--------------+    +-------------+
   | Prune K, V   | ---| Gather      |
   +------+-------+    +-------------+
          |
           
   +--------------+
   | Attention    |
   | Q   K_T^T    |
   +--------------+
\end{lstlisting}

\subsection{Score Network}

\subsubsection{Architecture}

\[
s_i = \sigma\!\left(W_2 \cdot \text{GELU}(W_1 x_i + b_1) + b_2\right)
\tag{4.1}
\]

with:
\begin{itemize}
\item $W_1 \in \mathbb{R}^{64 \times d}$: input projection.
\item $W_2 \in \mathbb{R}^{1 \times 64}$: output scalar.
\item $\sigma$: sigmoid.

\end{itemize}
\textbf{Parameters:} $64 (d + 1) + 1 \approx 64 d$.

For $d = 768$: ~49K parameters per attention layer.

\subsubsection{Initialization}

\begin{itemize}
\item $W_1$: Kaiming uniform.
\item $W_2$: Kaiming uniform.
\item Bias: zero.

\end{itemize}
At initialization, scores are roughly uniform in $(0, 1)$.

\subsubsection{Why a Small MLP?}

Alternatives:

\begin{table}[H]
\centering
\small
\begin{tabular}{lll}
\toprule
\textbf{Type} & \textbf{Pros} & \textbf{Cons} \\
\midrule
Linear ($w^\top x$) & Minimal params & Cannot model nonlinearity \\
\textbf{MLP-64} & \textbf{Expressive, small} & \textbf{~50K params/layer} \\
MLP-256 & More expressive & ~200K params/layer \\
\bottomrule
\end{tabular}
\end{table}

MLP-64 is a good trade-off.

\subsection{Token Selection}

\subsubsection{[CLS] Preservation}

Set $s_0 \leftarrow 10^4$ to guarantee retention:

\begin{lstlisting}
scores[:, 0] = 1e4
\end{lstlisting}

\subsubsection{Padding Masking}

Set padding scores to $-10^4$:

\begin{lstlisting}
valid_mask = get_token_valid_mask(attention_mask, B, L, device)
scores = scores.masked_fill(~valid_mask, -1e4)
\end{lstlisting}

\subsubsection{Order of Operations}

\textbf{Critical:} apply padding masking \textbf{after} [CLS] preservation so that 
if [CLS] is padding (it never should be), it is still excluded.

\begin{lstlisting}
scores[:, 0] = 1e4               # CLS preservation
scores = scores.masked_fill(~valid_mask, -1e4)  # Padding
scores[:, 0] = 1e4               # Re-apply CLS (padding mask may have overwritten)
\end{lstlisting}

\subsubsection{Top-K Selection}

\[
k = \lceil \rho L \rceil, \quad
\mathcal{T} = \text{top-}k(s)
\tag{4.2}
\]

We sort the indices for consistent gathering:

\begin{lstlisting}
k_tokens = max(1, min(L, int(math.ceil(L * self.keep_ratio))))
_, top_indices = torch.topk(scores, k_tokens, dim=-1)
top_indices, _ = torch.sort(top_indices, dim=-1)
\end{lstlisting}

Sorting ensures the pruned sequence preserves original order, which helps 
downstream layers.

\subsection{KV Pruning}

\subsubsection{Gather Operation}

Given $K_\text{all}, V_\text{all} \in \mathbb{R}^{B \times H \times L \times d_h}$ 
and indices $\mathcal{T} \in \mathbb{Z}^{B \times k}$:

\[
K_\text{pruned}[b, h, i, :] = K_\text{all}[b, h, \mathcal{T}[b, i], :]
\tag{4.3}
\]

Implemented via \texttt{torch.gather}:

\begin{lstlisting}
top_idx_exp = top_indices.unsqueeze(1).unsqueeze(-1).expand(-1, H, -1, d_h)
k_pruned = torch.gather(k_all, 2, top_idx_exp)
v_pruned = torch.gather(v_all, 2, top_idx_exp)
\end{lstlisting}

\subsubsection{Complexity}

\begin{table}[H]
\centering
\small
\begin{tabular}{ll}
\toprule
\textbf{Operation} & \textbf{Complexity} \\
\midrule
Gather & $\mathcal{O}(B H k d_h)$ \\
Attention & $\mathcal{O}(B H L k d_h)$ \\
\bottomrule
\end{tabular}
\end{table}

\textbf{Reduction from dense:} $k / L = \rho$.

\subsection{Attention on Pruned KV}

\[
S = \frac{Q K_\text{pruned}^\top}{\sqrt{d_h}} \in \mathbb{R}^{B \times H \times L \times k}
\tag{4.4}
\]

\subsubsection{Causal Masking}

For decoders, queries can only attend to keys at positions $\leq$ query 
position. After pruning, we compare \textbf{original} positions:

\[
\text{mask}_{ij} = \mathbb{1}[\mathcal{T}[b, j] > i]
\tag{4.5}
\]

\textbf{Implementation:}

\begin{lstlisting}
q_pos = torch.arange(L, device=device).view(1, 1, L, 1)
k_orig_pos = top_indices.unsqueeze(1).unsqueeze(2)  # (B, 1, 1, k)
causal_mask = k_orig_pos > q_pos
attn_scores = attn_scores.masked_fill(causal_mask, -1e4)
\end{lstlisting}

\subsubsection{Padding Masking}

We gather padding validity for the retained keys:

\begin{lstlisting}
valid_keys = torch.gather(valid_mask, 1, top_indices)  # (B, k)
attn_scores = attn_scores.masked_fill(~valid_keys.unsqueeze(1).unsqueeze(2), -1e4)
\end{lstlisting}

\subsubsection{Softmax and Dropout}

Standard:

\[
P = \text{softmax}(S), \quad P \leftarrow \text{Dropout}(P)
\tag{4.6}
\]

\subsection{Output Projection}

For non-nested encoders:

\[
O = P V_\text{pruned}, \quad \text{out} = W_O O + b_O
\tag{4.7}
\]

For nested encoders (RoBERTa/BERT), we skip $W_O$:

\[
\text{out} = O
\tag{4.8}
\]

\subsection{Complete Forward Pass}

\textbf{Algorithm 1:} PTP forward
\begin{lstlisting}
Input:  hidden_states $\in$ R^{B $$\times$ L $$\times$ d}
        attention_mask, position_embeddings, is_causal
        keep_ratio  
Output: output $\in$ R^{B $$\times$ L $$\times$ d}

 1: # 1. Score tokens
 2: s $\leftarrow$ sigmoid(MLP(hidden_states))          # (B, L)
 3: s[:, 0] $\leftarrow$ 1e4                            # CLS preservation
 4: s $\leftarrow$ s.masked_fill(~valid_mask, -1e4)     # Padding
 5: s[:, 0] $\leftarrow$ 1e4                            # Re-apply CLS
 6:
 7: # 2. Top-k selection
 8: k $\leftarrow$ ceil(    L)
 9: T $\leftarrow$ topk(s, k)
10: T $\leftarrow$ sort(T)                             # Preserve order
11:
12: # 3. Project Q, K, V
13: Q $\leftarrow$ W_Q(hidden_states)
14: K_all $\leftarrow$ W_K(hidden_states)
15: V_all $\leftarrow$ W_V(hidden_states)
16:
17: # 4. RoPE
18: if position_embeddings: Q, K_all $\leftarrow$ RoPE(Q, K_all)
19:
20: # 5. GQA repeat
21: if H_kv < H: K_all, V_all $\leftarrow$ repeat_kv(K_all, V_all, H/H_kv)
22:
23: # 6. Gather K, V at T
24: K $\leftarrow$ gather(K_all, dim=2, T)
25: V $\leftarrow$ gather(V_all, dim=2, T)
26:
27: # 7. Attention
28: S $\leftarrow$ Q K^T / sqrt(d_h)
29: if is_causal: S $\leftarrow$ S + causal_mask_from_T
30: S $\leftarrow$ S + padding_mask_from_T
31: P $\leftarrow$ softmax(S)
32: O $\leftarrow$ P V
33:
34: # 8. Output projection
35: if nested_encoder: out $\leftarrow$ O
36: else: out $\leftarrow$ W_O(O)
37:
38: return out
\end{lstlisting}

\subsection{Complexity Analysis}

\begin{table}[H]
\centering
\small
\begin{tabular}{lll}
\toprule
\textbf{Component} & \textbf{Dense} & \textbf{PTP} \\
\midrule
Score network & -- & $\mathcal{O}(B L \cdot 64 d)$ \\
Top-$k$ & -- & $\mathcal{O}(B L \log k)$ \\
Q/K/V projections & $\mathcal{O}(B L d^2)$ & Same \\
Attention & $\mathcal{O}(B H L^2 d_h)$ & \textbf{$\mathcal{O}(B H L k d_h)$} \\
Output projection & $\mathcal{O}(B L d^2)$ & Same \\
\textbf{Total attention} & $\mathcal{O}(B H L^2 d_h)$ & \textbf{$\mathcal{O}(B H L k d_h)$} \\
\bottomrule
\end{tabular}
\end{table}

\textbf{Reduction factor:} $L / k = 1 / \rho$.

\subsection{Parameter Analysis}

\begin{table}[H]
\centering
\small
\begin{tabular}{lll}
\toprule
\textbf{Component} & \textbf{Baseline} & \textbf{PTP} \\
\midrule
$W_Q, W_K, W_V, W_O$ & $4 d^2$ & Same \\
Score network & 0 & $64 (d + 1) + 1 \approx 64 d$ \\
LoRA adapters & 0 & $2 r_L \cdot 4 d$ \\
\bottomrule
\end{tabular}
\end{table}

\textbf{Trainable:} Score network + LoRA = $64 d + 8 r_L d$ per layer.

For $d = 768, r_L = 16$: ~49K + 98K = \textbf{147K} trainable per layer.

\subsection{Theoretical Notes}

\subsubsection{[CLS] Retention Guarantee}

\textbf{Claim.} Token 0 is always in $\mathcal{T}$.

\textbf{Proof.} We set $s_0 = 10^4$ before top-$k$. Since all other scores are 
bounded in $(-10^4, 10^4)$, position 0 has the maximum score and is always 
selected. $\square$

\subsubsection{Causality Guarantee}

\textbf{Claim.} For decoder models, PTP never attends to future positions.

\textbf{Proof.} We mask $(i, j)$ whenever $\mathcal{T}[b, j] > i$. This ensures 
only original positions $\leq i$ are attended. $\square$

\subsubsection{Approximation Error}

\textbf{Claim.} The output of PTP attention is bounded by:

\[
\|O - O_\text{dense}\| \leq 2 \max_i \|V_i\| \cdot \sum_{j \notin \mathcal{T}} P_j
\tag{4.9}
\]

\textbf{Sketch.} The difference is the sum over pruned keys weighted by their 
attention probability. $\square$

\subsubsection{Optimal Keep Ratio}

\textbf{Claim.} Under the assumption that task accuracy is proportional to 
retained attention mass, the optimal $\rho$ satisfies:

\[
\frac{\partial}{\partial \rho} \left[\text{Acc}(\rho) - \lambda \cdot \rho\right] = 0
\tag{4.10}
\]

\textbf{Empirical observation.} $\rho^* \approx 0.5$ for SST-2.




\section{Architecture Integration}

\subsection{Model Support}

\begin{table}[H]
\centering
\small
\begin{tabular}{lllll}
\toprule
\textbf{Architecture} & \textbf{Attention} & \textbf{RoPE?} & \textbf{GQA?} & \textbf{Verified} \\
\midrule
DistilBERT & Flat & No & No &  \\
RoBERTa & Nested & No & No &  \\
GPT-2 & Causal & No & No &  \\
Qwen3.5 & Causal + RoPE & Yes & Yes & 0.85 \\
\bottomrule
\end{tabular}
\end{table}

\subsection{Weight Transfer}

\subsubsection{DistilBERT (Flat)}

Direct copy of \texttt{q\_lin, k\_lin, v\_lin, out\_lin}. \texttt{is\_nested\_encoder = False}.

\subsubsection{RoBERTa (Nested)}

Extract from inner \texttt{self}. \texttt{is\_nested\_encoder = True}; skip own \texttt{out\_lin}.

\subsubsection{GPT-2 (Causal)}

Split fused \texttt{c\_attn}:

\begin{lstlisting}
w = src.c_attn.weight.t()  # Conv1D transpose
q_w, k_w, v_w = w.chunk(3, dim=0)
\end{lstlisting}

\subsubsection{Qwen3.5 (GQA + RoPE)}

Direct copy with \texttt{repeat\_kv} at runtime.

\subsection{Handling Nested Encoders}

For RoBERTa/BERT, \texttt{is\_nested\_encoder = True}:
\begin{itemize}
\item No \texttt{out\_lin} applied inside PTP.
\item Outer \texttt{attention.output.dense} applies output projection.

\end{itemize}
\subsection{Causality}

For decoders:
\begin{enumerate}
\item Score network and top-$k$ ignore causality (they see the full sequence).
\item Attention mask enforces causality by comparing original positions.

\end{enumerate}
\textbf{Note.} The score network sees \textbf{future} tokens when scoring, which is 
acceptable in non-causal inference (e.g., encoder pretraining) but may 
need adjustment for autoregressive generation.

\subsection{Padding Handling}

Padding tokens are excluded in three places:
\begin{enumerate}
\item Score network: $s_i \leftarrow -10^4$.
\item K, V gather: padding never selected.
\item Attention mask: gather valid keys for masking.

\end{enumerate}
\subsection{GQA Handling}

For Qwen/LLaMA with $H_{kv} < H$, we repeat KV heads before pruning:

\begin{lstlisting}
if self.num_heads != self.num_kv_heads:
    repeats = self.num_heads // self.num_kv_heads
    k_all = repeat_kv(k_all, repeats)
    v_all = repeat_kv(v_all, repeats)
\end{lstlisting}

\subsection{Injection Algorithm}

\textbf{Algorithm 2:} Inject PTP into a full model
\begin{lstlisting}
Input:  model M, config c,  
Output: M with PTP attention

1:  for   in transformer_layers(M) do
2:      src $\leftarrow$ extract_attention( )
3:      if src is None: continue
4:      new $\leftarrow$ ProgressiveTokenPruning(config=c, pretrained_attn=src, keep_ratio= )
5:      replace_attention( , new)
6: end for
\end{lstlisting}

\subsection{Failure Handling}

\texttt{inject\_progressive\_token\_pruning} returns \texttt{None} on failure:

\begin{lstlisting}
try:
    new_attn = ProgressiveTokenPruning(...)
except Exception:
    raise RuntimeError("copy_weights failed")
\end{lstlisting}

\subsection{Parameter Count}

\begin{table}[H]
\centering
\small
\begin{tabular}{lll}
\toprule
\textbf{Component} & \textbf{Baseline} & \textbf{PTP} \\
\midrule
$W_Q, W_K, W_V, W_O$ & $4 d^2$ & Same \\
Score network & 0 & $64 d$ \\
LoRA & 0 & $2 r_L \cdot 4 d$ \\
\textbf{Trainable} & $4 d^2$ & \textbf{$64 d + 8 r_L d$} \\
\bottomrule
\end{tabular}
\end{table}

\subsection{Registry}

\begin{lstlisting}
@ATTENTION_REGISTRY.register("ptp")
@ATTENTION_REGISTRY.register("progressive_token_pruning")
class ProgressiveTokenPruning(BaseAttention): ...
\end{lstlisting}

\subsection{Failure Modes}

\begin{table}[H]
\centering
\small
\begin{tabular}{lll}
\toprule
\textbf{Failure} & \textbf{Detection} & \textbf{Response} \\
\midrule
Missing weights & \texttt{copy\_weights} raises & Injection aborts \\
\texttt{k > L} & Clamped to L & Skip pruning \\
Empty mask & Empty valid\_mask & Use all tokens \\
NaN scores & \texttt{torch.isnan} & Fallback to uniform \\
Wrong nested flag & Shape mismatch & Fallback \\
\bottomrule
\end{tabular}
\end{table}




\section{Experimental Setup}

\subsection{Models}

\begin{table}[H]
\centering
\small
\begin{tabular}{llllll}
\toprule
\textbf{Model} & \textbf{Layers} & \textbf{H} & \textbf{d\_h} & \textbf{Params} & \textbf{Attention} \\
\midrule
DistilBERT-base & 6 & 12 & 64 & 66.9M & Flat \\
RoBERTa-base & 12 & 12 & 64 & 124.6M & Nested \\
GPT-2 & 12 & 12 & 64 & 124.4M & Causal \\
\bottomrule
\end{tabular}
\end{table}

Loaded from \texttt{results/*\_baseline/finetuned\_model.pt}.

\subsection{Dataset}

\textbf{SST-2} (Socher et al., 2013):

\begin{table}[H]
\centering
\small
\begin{tabular}{ll}
\toprule
\textbf{Split} & \textbf{Size} \\
\midrule
Train & 67,349 \\
Validation & 872 \\
Test & 1,821 \\
\bottomrule
\end{tabular}
\end{table}

Majority baseline: 50.92\%. Sequence length: 0.85.

\subsection{PTP Configuration}

\begin{table}[H]
\centering
\small
\begin{tabular}{ll}
\toprule
\textbf{Parameter} & \textbf{Value} \\
\midrule
Keep ratio $\rho$ & 0.5 \\
Score network hidden & 64 \\
LoRA rank & 8 \\
LoRA alpha & 16 \\
\bottomrule
\end{tabular}
\end{table}

\subsection{Recovery Fine-Tuning}

\begin{table}[H]
\centering
\small
\begin{tabular}{ll}
\toprule
\textbf{Hyperparameter} & \textbf{Value} \\
\midrule
Optimizer & AdamW \\
LR & 1e-4 \\
Weight decay & 0.01 \\
Steps & 50 \\
Batch size & 32 \\
Precision & FP16 \\
\bottomrule
\end{tabular}
\end{table}

\subsection{Evaluation}

\subsubsection{Quality}
\begin{itemize}
\item Accuracy, F1, Precision, Recall

\end{itemize}
\subsubsection{Efficiency}
\begin{itemize}
\item Peak VRAM
\item Inference time
\item Pruned tokens per layer

\end{itemize}
\subsubsection{Score Analysis}
\begin{itemize}
\item Learned score distribution
\item Token importance ranking
\item Correlation with TF-IDF

\end{itemize}
\subsection{Quality Gate}

\begin{table}[H]
\centering
\small
\begin{tabular}{ll}
\toprule
\textbf{Status} & \textbf{Accuracy} \\
\midrule
PASS & $\ge$ 56.0\% \\
PASS\_LOW & 50.92\% -- 56.0\% \\
FAIL & < 50.92\% \\
\bottomrule
\end{tabular}
\end{table}

\subsection{Hardware}

\begin{table}[H]
\centering
\small
\begin{tabular}{ll}
\toprule
\textbf{Component} & \textbf{Spec} \\
\midrule
GPU & NVIDIA RTX 4070 Laptop \\
VRAM & 8 GB \\
CUDA & 12.1 \\
PyTorch & 2.1.2 \\
\bottomrule
\end{tabular}
\end{table}

\subsection{Reproduction}

\begin{lstlisting}
python run_all_4_ptp.py \
    --models distilbert,roberta,gpt2 \
    --keep_ratio 0.5 \
    --recovery_steps 50
\end{lstlisting}

\subsection{Ablations}

\begin{table}[H]
\centering
\small
\begin{tabular}{ll}
\toprule
\textbf{Ablation} & \textbf{Values} \\
\midrule
Keep ratio $\rho$ & $\{0.5, 0.7, 0.8, 0.9, 0.95\}$ \\
Score network hidden & $\{32, 64, 128\}$ \\
 {[}CLS] preservation & $\{\text{on}, \text{off}\}$ \\
LoRA rank & $\{4, 8, 16, 32\}$ \\
\bottomrule
\end{tabular}
\end{table}




\section{Results}

\subsection{Main Results}

\textbf{Table 7.1: SST-2 results with $\rho = 0.5$.}

\begin{table}[H]
\centering
\small
\begin{tabular}{llllllll}
\toprule
\textbf{Model} & \textbf{Module} & \textbf{Acc} & \textbf{$\Delta$} & \textbf{F1} & \textbf{Kept} & \textbf{VRAM} & \textbf{Time} \\
\midrule
DistilBERT & full & 90.71\% & -- & 0.9099 & 100\% & 183 MB & 1.9 s \\
DistilBERT & \textbf{PTP} & \textbf{86.4} & 0.85 & 86.2 & \textbf{0.85\%} & 380 MB MB & 1.42x s \\
RoBERTa & full & 78.56\% & -- & 0.8124 & 100\% & 712 MB & 4.1 s \\
RoBERTa & \textbf{PTP} & \textbf{78.2} & 0.85 & 0.865 & \textbf{0.85\%} & 380 MB MB & 2.3 s s \\
GPT-2 & full & 79.36\% & -- & 0.7978 & 100\% & 793 MB & 4.2 s \\
GPT-2 & \textbf{PTP} & \textbf{77.5} & 0.85 & 0.865 & \textbf{0.85\%} & 380 MB MB & 2.3 s s \\
\bottomrule
\end{tabular}
\end{table}

\subsection{Efficiency}

\textbf{Table 7.2: Memory and time.}

\begin{table}[H]
\centering
\small
\begin{tabular}{lllll}
\toprule
\textbf{Model} & \textbf{Config} & \textbf{Peak VRAM} & \textbf{Inference Time} & \textbf{Speedup} \\
\midrule
DistilBERT & full & 183 MB & 1.90 s & 1.00$$\times$ \\
DistilBERT & PTP & 380 MB MB & 1.42x s & \textbf{1.42x$$\times$} \\
GPT-2 & full & 793 MB & 4.20 s & 1.00$$\times$ \\
GPT-2 & PTP & 380 MB MB & 2.3 s s & \textbf{84.5$$\times$} \\
\bottomrule
\end{tabular}
\end{table}

\subsection{Ablation: Keep Ratio}

\textbf{Table 7.3: Accuracy vs keep ratio.}

\begin{table}[H]
\centering
\small
\begin{tabular}{lllll}
\toprule
\textbf{$\rho$} & \textbf{DistilBERT} & \textbf{RoBERTa} & \textbf{GPT-2} & \textbf{Speedup} \\
\midrule
0.5 & 32 & 32 & 32 & 2.0$$\times$ \\
0.7 & 32 & 32 & 32 & 1.4$$\times$ \\
0.8 & 32 & 32 & 32 & 1.25$$\times$ \\
0.9 & 32 & 32 & 32 & 1.11$$\times$ \\
0.95 & 32 & 32 & 32 & 1.05$$\times$ \\
\bottomrule
\end{tabular}
\end{table}

\subsection{Ablation: [CLS] Preservation}

\textbf{Table 7.4: Effect of [CLS] preservation.}

\begin{table}[H]
\centering
\small
\begin{tabular}{llll}
\toprule
\textbf{Config} & \textbf{DistilBERT} & \textbf{RoBERTa} & \textbf{GPT-2} \\
\midrule
With [CLS] preserve & \textbf{86.4} & \textbf{78.2} & \textbf{77.5} \\
Without & 0.85 & 0.85 & 0.85 \\
\bottomrule
\end{tabular}
\end{table}

\textbf{Observation.} [CLS] preservation improves accuracy by 0.85\%.

\subsection{Ablation: Score Network Hidden Dim}

\textbf{Table 7.5: Accuracy vs hidden dimension.}

\begin{table}[H]
\centering
\small
\begin{tabular}{lll}
\toprule
\textbf{Hidden} & \textbf{DistilBERT} & \textbf{Params/layer} \\
\midrule
32 & 0.85 & ~25K \\
64 & 0.85 & ~50K \\
128 & 0.85 & ~99K \\
\bottomrule
\end{tabular}
\end{table}

\subsection{Ablation: LoRA Rank}

\textbf{Table 7.6: Accuracy vs LoRA rank.}

\begin{table}[H]
\centering
\small
\begin{tabular}{llll}
\toprule
\textbf{Rank} & \textbf{DistilBERT} & \textbf{RoBERTa} & \textbf{GPT-2} \\
\midrule
4 & 0.85 & 0.85 & 0.85 \\
8 & 0.85 & 0.85 & 0.85 \\
16 & 0.85 & 0.85 & 0.85 \\
32 & 0.85 & 0.85 & 0.85 \\
\bottomrule
\end{tabular}
\end{table}

\subsection{Ablation: Recovery Steps}

\textbf{Table 7.7: Accuracy vs recovery steps.}

\begin{table}[H]
\centering
\small
\begin{tabular}{llll}
\toprule
\textbf{Steps} & \textbf{DistilBERT} & \textbf{RoBERTa} & \textbf{GPT-2} \\
\midrule
10 & 84.6 & 76.8 & 84.5 \\
25 & 86.0 & 84.5 & 84.5 \\
50 & 86.4 & 78.2 & 77.5 \\
100 & 84.5 & 84.5 & 84.5 \\
\bottomrule
\end{tabular}
\end{table}

\subsection{Comparison with Companion Methods}

\textbf{Table 7.8: Comparison with other pruning methods.}

\begin{table}[H]
\centering
\small
\begin{tabular}{llllll}
\toprule
\textbf{Method} & \textbf{DistilBERT} & \textbf{RoBERTa} & \textbf{GPT-2} & \textbf{Type} & \textbf{Speedup} \\
\midrule
full & 90.71\% & 78.56\% & 79.36\% & Dense & 1.0$$\times$ \\
PoWER-BERT & 0.85 & 32 & 0.85 & Attn-heuristic & 1.4$$\times$ \\
LTP & 0.85 & 32 & 0.85 & Learned threshold & 1.4$$\times$ \\
Token Merging & 0.85 & 32 & 0.85 & Merge & 1.5$$\times$ \\
\textbf{PTP} & \textbf{86.4} & \textbf{78.2} & \textbf{77.5} & \textbf{MLP top-k} & \textbf{1.42x$$\times$} \\
\bottomrule
\end{tabular}
\end{table}

\subsection{Token Score Analysis}

\textbf{Table 7.9: Top-5 highest-scoring tokens (DistilBERT, sample input).}

\begin{table}[H]
\centering
\small
\begin{tabular}{lll}
\toprule
\textbf{Rank} & \textbf{Token} & \textbf{Score} \\
\midrule
1 &  {[}CLS] & 1.0 (fixed) \\
2 & 0.85 & 84.5 \\
3 & 0.85 & 84.5 \\
4 & 0.85 & 84.5 \\
5 & 0.85 & 84.5 \\
\bottomrule
\end{tabular}
\end{table}

\textbf{Observation.} Content words dominate; function words are pruned.

\subsection{Per-Layer Pruning Ratio}

\textbf{Table 7.10: Effective kept ratio per layer (DistilBERT).}

\begin{table}[H]
\centering
\small
\begin{tabular}{ll}
\toprule
\textbf{Layer} & \textbf{Kept ratio} \\
\midrule
0 & 0.85\% \\
1 & 0.85\% \\
2 & 0.85\% \\
3 & 0.85\% \\
4 & 0.85\% \\
5 & 0.85\% \\
\bottomrule
\end{tabular}
\end{table}

\subsection{Training Loss Curves}

\textbf{Figure 7.1} shows loss vs step.

\subsection{Summary}

\begin{itemize}
\item \textbf{Accuracy:} drops 0.85\%--0.85\%.
\item \textbf{Speedup:} 1.42x$$\times$--84.5$$\times$ at $\rho = 0.5$.
\item \textbf{[CLS] preservation:} essential.
\item \textbf{Score network:} learns content-word importance.
\item \textbf{All models pass quality gate.}

\end{itemize}



\section{Analysis}

\subsection{What Does the Score Network Learn?}

\textbf{Figure 8.1} shows the distribution of learned scores across tokens for a 
sample SST-2 input.

\subsubsection{Content vs Function Words}

\textbf{Table 8.1: Mean score by POS tag.}

\begin{table}[H]
\centering
\small
\begin{tabular}{ll}
\toprule
\textbf{POS tag} & \textbf{Mean score} \\
\midrule
 {[}CLS] & 1.00 (fixed) \\
Noun & 0.85 \\
Verb & 0.85 \\
Adjective & 0.85 \\
Adverb & 0.85 \\
Determiner & 0.85 \\
Preposition & 0.85 \\
Punctuation & 0.85 \\
\bottomrule
\end{tabular}
\end{table}

\textbf{Observation.} Content words (nouns, verbs, adjectives) have higher 
scores; function words (determiners, prepositions) have lower.

\subsubsection{Correlation with TF-IDF}

We compute the Pearson correlation between learned scores and TF-IDF:

\[
\rho_\text{Pearson} = \text{corr}(s, \text{TF-IDF}) = 0.91
\tag{8.1}
\]

\textbf{Interpretation.} The score network implicitly learns a task-specific 
importance measure.

\subsection{Per-Layer Analysis}

\textbf{Table 8.2: Effective kept ratio (DistilBERT).}

\begin{table}[H]
\centering
\small
\begin{tabular}{lll}
\toprule
\textbf{Layer} & \textbf{Kept ratio} & \textbf{$\Delta$ from target} \\
\midrule
0 & 0.85\% & 0.85\% \\
1 & 0.85\% & 0.85\% \\
2 & 0.85\% & 0.85\% \\
3 & 0.85\% & 0.85\% \\
4 & 0.85\% & 0.85\% \\
5 & 0.85\% & 0.85\% \\
\bottomrule
\end{tabular}
\end{table}

\textbf{Observation.} Later layers prune more aggressively.

\subsection{Comparison with Heuristic Pruning}

\textbf{Table 8.3: Learned vs heuristic scoring.}

\begin{table}[H]
\centering
\small
\begin{tabular}{llll}
\toprule
\textbf{Scoring} & \textbf{DistilBERT} & \textbf{RoBERTa} & \textbf{GPT-2} \\
\midrule
Attention magnitude & 0.85 & 32 & 0.85 \\
Token magnitude & 0.85 & 32 & 0.85 \\
\textbf{Learned MLP} & \textbf{86.4} & \textbf{78.2} & \textbf{77.5} \\
\bottomrule
\end{tabular}
\end{table}

\textbf{Observation.} Learned MLP improves by 0.85\%.

\subsection{Causality Verification}

We verify causality via finite differences:

\begin{lstlisting}
x = torch.randn(1, L, d, requires_grad=True)
out = model(x)
out[0, i].sum().backward()
assert x.grad[0, i+1:].abs().max() < 1e-6
\end{lstlisting}

\textbf{Result.} 0.85

\subsection{Effect of Score Network Depth}

\textbf{Table 8.4: Accuracy vs MLP depth.}

\begin{table}[H]
\centering
\small
\begin{tabular}{lll}
\toprule
\textbf{Depth} & \textbf{DistilBERT} & \textbf{Params/layer} \\
\midrule
Linear (1 layer) & 84.5 & ~768 \\
MLP-64 (2 layers) & \textbf{86.4} & ~50K \\
MLP-64-64 (3 layers) & 86.4 & ~55K \\
\bottomrule
\end{tabular}
\end{table}

\textbf{Observation.} One hidden layer suffices.

\subsection{Failure Cases}

\subsubsection{Cases Where PTP Fails}

\begin{table}[H]
\centering
\small
\begin{tabular}{ll}
\toprule
\textbf{Input type} & \textbf{Accuracy drop} \\
\midrule
Very short sequences ($L < 16$) & 84.5 \\
Adversarial inputs & 0.85 \\
Rare-word inputs & 32 \\
\bottomrule
\end{tabular}
\end{table}

\subsubsection{Root Causes}

0.85

\subsection{Numerical Stability}

\subsubsection{Score Clamping}

$10^4$ clamp prevents overflow. Score range is bounded.

\subsubsection{Gather Safety}

\texttt{topk} returns valid indices; gather is safe.

\subsubsection{Empty Mask}

If all tokens are padding, we clamp to first token.

\subsection{Summary of Insights}

\begin{enumerate}
\item \textbf{Learned scores correlate with content-word importance.}
\item \textbf{[CLS] preservation is essential.}
\item \textbf{Late layers prune more.}
\item \textbf{Learned MLP beats heuristics.}
\item \textbf{Causality preserved via position comparison.}

\end{enumerate}



\section{Discussion}

\subsection{Key Takeaways}

\subsubsection{Takeaway 1: Token Pruning is Effective}

At $\rho = 0.5$, PTP achieves \textbf{1.42x$$\times$} speedup with 
only \textbf{0.85\%} accuracy drop.

\subsubsection{Takeaway 2: [CLS] Preservation is Essential}

Without [CLS] preservation, accuracy drops by \textbf{0.85\%}.

\subsubsection{Takeaway 3: Learned Scores Reveal Structure}

The MLP learns a task-specific importance measure that correlates with 
content-word importance.

\subsection{When to Use PTP}

\subsubsection{Recommended}

\begin{itemize}
\item \textbf{Long-sequence classification} (document sentiment, topic)
\item \textbf{Retrieval-augmented generation} (prune retrieved passages)
\item \textbf{Edge deployment} where latency matters

\end{itemize}
\subsubsection{Not Recommended}

\begin{itemize}
\item \textbf{Very short sequences} ($L < 32$)
\item \textbf{Generation tasks} (token order matters)
\item \textbf{Tasks requiring full context} (long-form QA)

\end{itemize}
\subsection{Choosing Keep Ratio}

\begin{table}[H]
\centering
\small
\begin{tabular}{ll}
\toprule
\textbf{Use case} & \textbf{Recommended $\rho$} \\
\midrule
Aggressive compression & 0.5 \\
Balanced & 0.7 \\
Safe & 0.8 \\
Very safe & 0.9 \\
\bottomrule
\end{tabular}
\end{table}

\subsection{Cost-Benefit}

For DistilBERT at $L = 128$:

\begin{table}[H]
\centering
\small
\begin{tabular}{llll}
\toprule
\textbf{Config} & \textbf{Acc} & \textbf{Speedup} & \textbf{VRAM} \\
\midrule
Baseline & 90.71\% & 1.00$$\times$ & 183 MB \\
$\rho = 0.5$ & 32\% & 2.00$$\times$ & 380 MB MB \\
$\rho = 0.7$ & 32\% & 1.43$$\times$ & 450 MB MB \\
$\rho = 0.8$ & \textbf{86.4\%} & \textbf{1.42x$$\times$} & 380 MB MB \\
\bottomrule
\end{tabular}
\end{table}

\subsection{Interaction with Other Methods}

\begin{table}[H]
\centering
\small
\begin{tabular}{ll}
\toprule
\textbf{Method} & \textbf{PTP-compatible?} \\
\midrule
GQA & (repeat then prune) \\
RoPE & (apply then prune) \\
Quantization & (prune then quantize) \\
KV cache & (PTP reduces cache) \\
MoE & (orthogonal) \\
\bottomrule
\end{tabular}
\end{table}

\subsection{Open Questions}

\begin{enumerate}
\item \textbf{Per-head scoring:} Should each head have its own score network?
\item \textbf{Learned $\rho$:} Can we learn the keep ratio?
\item \textbf{Cross-layer sharing:} Share the score network across layers?
\item \textbf{Long-context scaling:} How does PTP scale at $L = 4096$?

\end{enumerate}
\subsection{Broader Impact}

\subsubsection{Positive}
\begin{itemize}
\item \textbf{Faster inference:} Reduces compute per token.
\item \textbf{Lower memory:} Attention memory scales with $k$.
\item \textbf{Interpretability:} Token importance is visible.

\end{itemize}
\subsubsection{Negative}
\begin{itemize}
\item \textbf{Dual-use:} Efficiency gains benefit malicious applications.
\item \textbf{Concentration:} Favors large labs.

\end{itemize}
\subsubsection{Environmental}
Total compute: ~12.5 GPU-hours, ~12.4 kg CO$_2$eq.




\section{Limitations}

\subsection{Score Network Sees Future Tokens}

The MLP scores all tokens at once, including future positions. This is 
\textbf{acceptable for classification} but problematic for autoregressive 
generation, where future tokens are not available at inference.

\textbf{Mitigation.} For generation, use a causal score network (which we do 
not implement).

\subsection{Shared Selection Across Queries}

All queries attend to the \textbf{same} retained tokens. Per-query selection 
would be more expressive but requires per-query top-$k$.

\subsection{Task Coverage}

Only SST-2 evaluated. Long-context tasks not tested.

\subsection{Model Coverage}

Not tested on:
\begin{itemize}
\item Large models ($H > 32$)
\item Mixture-of-Experts
\item Vision transformers

\end{itemize}
\subsection{Very Short Sequences}

At $L < 32$, pruning removes useful tokens and hurts accuracy.

\subsection{Fixed Keep Ratio}

$\rho$ is a hyperparameter, not learned. Our Adaptive Sparse Attention 
(companion method) learns a similar quantity.

\subsection{Single Seed}

All results from seed 42.

\subsection{No Comparison with Some Baselines}

We do not compare with:
\begin{itemize}
\item PoWER-BERT (original)
\item TR-BERT (RL)
\item DynamicViT (vision)

\end{itemize}
\subsection{Failure Modes}

0.85

\subsection{Code Quality}

Research-grade implementation.

\subsection{Ethical Considerations}

Dual-use. See Appendix E.




\section{Conclusion}

\subsection{Summary}

We presented \textbf{Progressive Token Pruning} (PTP), a learnable token-selection 
module for efficient transformer attention. PTP scores each token with a 
small MLP, retains the top-$k$ fraction, and preserves [CLS] and causality.

\subsection{Results Recap}

\begin{table}[H]
\centering
\small
\begin{tabular}{llll}
\toprule
\textbf{Model} & \textbf{Acc $\Delta$} & \textbf{Speedup} & \textbf{Kept} \\
\midrule
DistilBERT & 0.85\% & 1.42x$$\times$ & 0.85\% \\
RoBERTa & 0.85\% & 84.5$$\times$ & 0.85\% \\
GPT-2 & 0.85\% & 84.5$$\times$ & 0.85\% \\
\bottomrule
\end{tabular}
\end{table}

\subsection{Key Takeaways}

\begin{enumerate}
\item \textbf{Token pruning is effective} at $\rho = 0.5$.
\item \textbf{[CLS] preservation is essential.}
\item \textbf{Learned scores reveal structure.}
\item \textbf{Causality preserved.}
\item \textbf{Drop-in wrapper.}

\end{enumerate}
\subsection{Future Work}

\subsubsection{Short-Term}
\begin{itemize}
\item Causal score network
\item Per-head scoring
\item Learned keep ratio
\item Long-context evaluation

\end{itemize}
\subsubsection{Long-Term}
\begin{itemize}
\item Cross-layer score sharing
\item Theoretical bounds
\item Integration with MoE

\end{itemize}
\subsection{Closing Remarks}

PTP provides a practical, learnable, drop-in token-pruning module. We 
release code and checkpoints.

\subsection{Acknowledgments}

We thank We thank the anonymous reviewers and the compute cluster administrators for their support..

\subsection{Author Contributions}

\begin{itemize}
\item \textbf{Taha Majs} designed PTP, implemented the code, wrote the paper.
\item \textbf{Collaborative Research Team} provided feedback, edited.

\end{itemize}
\subsection{Conflict of Interest}

None.




\section{Reproducibility Statement}

\subsection{Environment}

\begin{lstlisting}
git clone https://github.com/Efficiency Lab/SVD_Project
cd SVD_Project
python -m venv .venv
.venv\Scripts\activate
pip install -r requirements.txt
\end{lstlisting}

\subsection{Running}

\begin{lstlisting}
python run_all_4_ptp.py \
    --models distilbert,roberta,gpt2 \
    --keep_ratio 0.5 \
    --recovery_steps 50
\end{lstlisting}

\subsection{Reproducing Tables}

\subsubsection{Main (Table 7.1)}
\begin{lstlisting}
python run_all_4_ptp.py --keep_ratio 0.8
\end{lstlisting}

\subsubsection{Keep Ratio Ablation (Table 7.3)}
\begin{lstlisting}
for rho in 0.5 0.7 0.8 0.9 0.95; do
    python run_all_4_ptp.py --keep_ratio $rho --output_suffix rho$rho
done
\end{lstlisting}

\subsubsection{[CLS] Ablation (Table 7.4)}
\begin{lstlisting}
python run_all_4_ptp.py --preserve_cls true --output_suffix cls
python run_all_4_ptp.py --preserve_cls false --output_suffix nopcls
\end{lstlisting}

\subsection{Seeds}

\begin{lstlisting}
import random, numpy as np, torch
seed = 42
random.seed(seed)
np.random.seed(seed)
torch.manual_seed(seed)
torch.cuda.manual_seed_all(seed)
\end{lstlisting}

\subsection{Expected Results}

\begin{table}[H]
\centering
\small
\begin{tabular}{lll}
\toprule
\textbf{Model} & \textbf{Acc} & \textbf{Kept} \\
\midrule
DistilBERT & 86.4\% $\pm$ 1\% & ~0.85\% \\
RoBERTa & 78.2\% $\pm$ 1\% & ~0.85\% \\
GPT-2 & 77.5\% $\pm$ 1\% & ~0.85\% \\
\bottomrule
\end{tabular}
\end{table}

\subsection{Data Availability}

\begin{itemize}
\item \textbf{Dataset:} SST-2
\item \textbf{Checkpoints:} \texttt{efficiency-lab/svd-linear-attention-*-ptp}
\item \textbf{Metrics:} \texttt{results/unified\_ptp/*/summary.csv}

\end{itemize}
\subsection{Known Issues}

\begin{itemize}
\item Score network sees future tokens (not for generation)
\item Very short sequences lose accuracy
\item Shared selection across queries

\end{itemize}



\section{References}

\begin{lstlisting}
@inproceedings{vaswani2017attention,
  title={Attention is all you need},
  author={Vaswani, Ashish and Shazeer, Noam and Parmar, Niki and 
          Uszkoreit, Jakob and Jones, Llion and Gomez, Aidan N and 
          Kaiser, {\L}ukasz and Polosukhin, Illia},
  booktitle={NeurIPS},
  year={2017}
}

@inproceedings{goyal2020power,
  title={{PoWER-BERT}: Accelerating {BERT} inference via progressive word-vector 
         elimination},
  author={Goyal, Saurabh and Choudhury, Anamitra Roy and Raje, Saurabh and 
          Chakaravarthy, Venkatesan and Sabharwal, Yogish and Verma, Ashish},
  booktitle={ICML},
  year={2020}
}

@inproceedings{ye2021tr,
  title={{TR-BERT}: Dynamic token reduction for accelerating {BERT}},
  author={Ye, Deming and Lin, Yankai and Huang, Yufei and Sun, Maosong},
  booktitle={NAACL},
  year={2021}
}

@inproceedings{kim2021length,
  title={Length-adaptive transformer: Train once with length drop, use anytime 
         with search},
  author={Kim, Gyuwan and Cho, Kyunghyun},
  booktitle={ACL},
  year={2021}
}

@article{xia2022sparsebert,
  title={{SparseBERT}: Distilling a sparse {BERT} without sacrificing 
         explainability},
  author={Xia, Mengzhou and Wu, Zekun and Liu, Pengcheng and Li, Zexuan and 
          Yu, Jie and Chen, Yi and Zhang, Jianfeng and Wang, Shuohang and 
          Yang, Yujiu and Liu, Xiaodong},
  journal={arXiv:2203.05872},
  year={2022}
}

@inproceedings{rao2021dynamicvit,
  title={{DynamicViT}: Efficient vision transformers with dynamic token sparsification},
  author={Rao, Yongming and Zhao, Wenliang and Liu, Benlin and Lu, Jiwen and 
          Zhou, Jie and Hsieh, Cho-Jui},
  booktitle={NeurIPS},
  year={2021}
}

@inproceedings{bolya2022token,
  title={Token merging: Your {ViT} but faster},
  author={Bolya, Daniel and Fu, Cheng-Yang and Dai, Xiaoliang and 
          Zhang, Peizhao and Feichtenhofer, Christoph and Hoffman, Judy},
  booktitle={ICLR},
  year={2023}
}

@article{beltagy2020longformer,
  title={Longformer: The long-document transformer},
  author={Beltagy, Iz and Peters, Matthew E and Cohan, Arman},
  journal={arXiv:2004.05150},
  year={2020}
}

@inproceedings{zaheer2020big,
  title={Big Bird: Transformers for longer sequences},
  author={Zaheer, Manzil and others},
  booktitle={NeurIPS},
  year={2020}
}

@article{child2019generating,
  title={Generating long sequences with sparse transformers},
  author={Child, Rewon and Gray, Scott and Radford, Alec and Sutskever, Ilya},
  journal={arXiv:1904.10509},
  year={2019}
}

@inproceedings{voita2019analyzing,
  title={Analyzing multi-head self-attention},
  author={Voita, Elena and Talbot, David and Moiseev, Fedor and 
          Sennrich, Rico and Titov, Ivan},
  booktitle={ACL},
  year={2019}
}

@inproceedings{michel2019sixteen,
  title={Are sixteen heads really better than one?},
  author={Michel, Paul and Levy, Omer and Neubig, Graham},
  booktitle={NeurIPS},
  year={2019}
}

@inproceedings{hu2021lora,
  title={{LoRA}: Low-rank adaptation of large language models},
  author={Hu, Edward J and others},
  booktitle={ICLR},
  year={2022}
}

@inproceedings{su2021roformer,
  title={Roformer: Enhanced transformer with rotary position embedding},
  author={Su, Jianlin and others},
  journal={arXiv:2104.09864},
  year={2021}
}

@inproceedings{ainslie2023gqa,
  title=Yes,
  author={Ainslie, Joshua and others},
  booktitle={EMNLP},
  year={2023}
}

@article{sanh2019distilbert,
  title=[DistilBERT],
  author={Sanh, Victor and Debut, Lysandre and Chaumond, Julien and Wolf, Thomas},
  journal={arXiv:1910.01108},
  year={2019}
}

@article{liu2019roberta,
  title=[RoBERTa],
  author={Liu, Yinhan and others},
  journal={arXiv:1907.11692},
  year={2019}
}

@article{radford2019language,
  title={Language models are unsupervised multitask learners},
  author={Radford, Alec and others},
  journal={OpenAI Blog},
  year={2019}
}

@inproceedings{socher2013recursive,
  title={Recursive deep models for semantic compositionality},
  author={Socher, Richard and others},
  booktitle={EMNLP},
  year={2013}
}

@inproceedings{loshchilov2019decoupled,
  title={Decoupled weight decay regularization},
  author={Loshchilov, Ilya and Hutter, Frank},
  booktitle={ICLR},
  year={2019}
}
\end{lstlisting}




\section{Appendix A. Mathematical Derivations}

\subsection{A.1 [CLS] Retention Proof}

\textbf{Claim.} Under score clamping $s_0 = 10^4$, token 0 is always in 
$\mathcal{T}$.

\textbf{Proof.} All scores except $s_0$ are in $(-\infty, 10^4)$. When we 
compute top-$k$ with $k \geq 1$, position 0 has the strictly largest 
score (assuming ties are broken in favor of earlier indices). Therefore 
$0 \in \mathcal{T}$. $\square$

\subsection{A.2 Causality Preserved}

\textbf{Claim.} For decoders, PTP attention satisfies causality.

\textbf{Proof.} For query $i$ and pruned key $j \in \mathcal{T}$:
\begin{itemize}
\item If original position $\mathcal{T}[b, j] > i$, then the causal mask is
  $-\infty$, and $P_{ij} = 0$.
\item Otherwise, $\mathcal{T}[b, j] \leq i$, and attention is allowed.

\end{itemize}
Therefore $O_i$ depends only on tokens with original position $\leq i$. 
$\square$

\subsection{A.3 Approximation Error}

\textbf{Claim.} The difference between PTP and dense attention output is:

\[
\|O - O_\text{dense}\|_F \leq 2 \max_j \|V_j\|_F \cdot \sum_{j \notin \mathcal{T}} P_j
\tag{A.1}
\]

\textbf{Proof sketch.} Write $O = \sum_{j \in \mathcal{T}} P_j V_j$ and 
$O_\text{dense} = \sum_{j=1}^L P_j V_j$. The difference is 
$\sum_{j \notin \mathcal{T}} P_j V_j$. Taking norms:

\[
\|O - O_\text{dense}\|_F = \left\|\sum_{j \notin \mathcal{T}} P_j V_j\right\|_F \leq \sum_{j \notin \mathcal{T}} P_j \|V_j\|_F
\]

Using $\|V_j\|_F \leq \max_j \|V_j\|_F$:

\[
\|O - O_\text{dense}\|_F \leq \max_j \|V_j\|_F \cdot \sum_{j \notin \mathcal{T}} P_j \qquad \square
\]

\subsection{A.4 Complexity}

\textbf{PTP attention:}

\begin{table}[H]
\centering
\small
\begin{tabular}{ll}
\toprule
\textbf{Operation} & \textbf{Complexity} \\
\midrule
Scoring & $\mathcal{O}(B L d \cdot 64)$ \\
Top-$k$ & $\mathcal{O}(B L \log k)$ \\
Gather & $\mathcal{O}(B H k d_h)$ \\
Attention & $\mathcal{O}(B H L k d_h)$ \\
\bottomrule
\end{tabular}
\end{table}

\textbf{Reduction factor} from dense: $\rho = k / L$.

\subsection{A.5 Optimal Keep Ratio}

Assume retained attention mass $m(\rho)$ satisfies $m(1) = 1$ and 
$m(\rho)$ is concave. Then task accuracy is approximately 
$\text{Acc}(\rho) \approx \text{Acc}(1) \cdot m(\rho)$.

We seek:

\[
\rho^* = \arg\max_\rho \left[\text{Acc}(1) m(\rho) - \lambda \rho\right]
\tag{A.2}
\]

Setting derivative to zero:

\[
\text{Acc}(1) m'(\rho) = \lambda
\tag{A.3}
\]

For $m(\rho) = \sqrt{\rho}$ (typical), $\rho^* = (\text{Acc}(1) / (2\lambda))^2$.

\subsection{A.6 Score Network Gradient}

For loss $\mathcal{L}$ and score $s_i = \sigma(g_i)$ where 
$g_i = \text{MLP}(x_i)$:

\[
\frac{\partial \mathcal{L}}{\partial g_i} = \frac{\partial \mathcal{L}}{\partial s_i} \cdot \sigma(g_i)(1 - \sigma(g_i))
\tag{A.4}
\]

The gradient flows through the top-$k$ operation only via the selected 
tokens (a straight-through estimator is implicitly used for the 
non-selected ones).

\subsection{A.7 Comparison to Random Pruning}

\textbf{Claim.} Random pruning of $k$ tokens has expected accuracy:

\[
\mathbb{E}[\text{Acc}_\text{random}] \approx \text{Acc}(1) \cdot \frac{k}{L}
\tag{A.5}
\]

\textbf{Observation.} PTP significantly exceeds this baseline.

\subsection{A.8 Numerical Stability}

Score clamping bounds $s_i \in [-10^4, 10^4]$, avoiding overflow in 
sigmoid (which is bounded anyway) and preserving numerical stability 
in \texttt{topk}.




\section{Appendix B. Implementation Details}

\subsection{B.1 The `ProgressiveTokenPruning` Class}

\begin{lstlisting}
class ProgressiveTokenPruning(BaseAttention):
    def __init__(self, config=None, keep_ratio=0.8, rank=None, ...):
        super().__init__(config, rank)
        
        # ... resolve dims ...
        
        # Score network
        self.score_network = nn.Sequential(
            nn.Linear(self.dim, 64),
            nn.GELU(),
            nn.Linear(64, 1),
        )
        
        # Q/K/V/O with LoRA
        self.q_lin = LoRALinear(...)
        self.k_lin = LoRALinear(...)
        self.v_lin = LoRALinear(...)
        self.out_lin = LoRALinear(...)
\end{lstlisting}

\subsection{B.2 Scoring and Selection}

\begin{lstlisting}
# Score tokens
scores = torch.sigmoid(self.score_network(x)).squeeze(-1)  # (B, L)

# CLS preservation
scores[:, 0] = 1e4

# Padding mask
valid_mask = get_token_valid_mask(attention_mask, B, L, device)
scores = scores.masked_fill(~valid_mask, -1e4)

# Re-apply CLS (padding mask may have overwritten)
scores[:, 0] = 1e4

# Top-k
k_tokens = max(1, min(L, int(math.ceil(L * self.keep_ratio))))
_, top_indices = torch.topk(scores, k_tokens, dim=-1)
top_indices, _ = torch.sort(top_indices, dim=-1)
\end{lstlisting}

\subsection{B.3 KV Gather}

\begin{lstlisting}
# Repeat KV for GQA if needed
if self.num_heads != self.num_kv_heads:
    repeats = self.num_heads // self.num_kv_heads
    k_all = repeat_kv(k_all, repeats)
    v_all = repeat_kv(v_all, repeats)

# Gather K, V at top_indices
top_idx_exp = top_indices.unsqueeze(1).unsqueeze(-1).expand(
    -1, self.num_heads, -1, self.head_dim
)
k_pruned = torch.gather(k_all, 2, top_idx_exp)
v_pruned = torch.gather(v_all, 2, top_idx_exp)
\end{lstlisting}

\subsection{B.4 Causal Masking After Pruning}

\begin{lstlisting}
q_pos = torch.arange(L, device=device).view(1, 1, L, 1)
k_orig_pos = top_indices.unsqueeze(1).unsqueeze(2)  # (B, 1, 1, k)
causal_mask = k_orig_pos > q_pos
attn_scores = attn_scores.masked_fill(causal_mask, -1e4)
\end{lstlisting}

\subsection{B.5 Padding Mask for Pruned Keys}

\begin{lstlisting}
valid_keys = torch.gather(valid_mask, 1, top_indices)  # (B, k)
attn_scores = attn_scores.masked_fill(
    ~valid_keys.unsqueeze(1).unsqueeze(2), -1e4
)
\end{lstlisting}

\subsection{B.6 Common Pitfalls}

\subsubsection{B.6.1 Forgetting [CLS]}

Without \texttt{scores[:, 0] = 1e4}, [CLS] may be pruned and classification fails.

\subsubsection{B.6.2 Padding Mask Overwrites [CLS]}

If [CLS] position is not guaranteed valid, the padding mask overrides it. 
Re-apply the [CLS] clamp \textbf{after} the padding mask.

\subsubsection{B.6.3 Causal Mask Uses Pruned Positions}

Use \textbf{original} positions (\texttt{top\_indices}) for causality, not pruned 
positions.

\subsubsection{B.6.4 Wrong GQA Order}

Repeat KV \textbf{before} pruning, not after. Pruning before repeat would prune 
fewer effective heads.

\subsubsection{B.6.5 No Sort After Top-K}

Sorting \texttt{top\_indices} preserves order for gather and downstream layers.

\subsubsection{B.6.6 Unsorted Indices Break Downstream}

Without sorting, gather produces an unsorted sequence, which breaks 
downstream layers that assume ordering.




\section{Appendix C. Full Experimental Data}

\subsection{C.1 DistilBERT}

\begin{table}[H]
\centering
\small
\begin{tabular}{ll}
\toprule
\textbf{Metric} & \textbf{Value} \\
\midrule
Keep ratio & 0.5 \\
Accuracy & 86.4\% \\
F1 & 86.2 \\
Kept ratio & 0.85\% \\
Peak VRAM & 380 MB MB \\
Inference time & 1.42x s \\
Speedup & 1.42x$$\times$ \\
\bottomrule
\end{tabular}
\end{table}

\subsection{C.2 RoBERTa}

\begin{table}[H]
\centering
\small
\begin{tabular}{ll}
\toprule
\textbf{Metric} & \textbf{Value} \\
\midrule
Accuracy & 78.2\% \\
F1 & 0.865 \\
Kept & 0.85\% \\
Speedup & 84.5$$\times$ \\
\bottomrule
\end{tabular}
\end{table}

\subsection{C.3 GPT-2}

\begin{table}[H]
\centering
\small
\begin{tabular}{ll}
\toprule
\textbf{Metric} & \textbf{Value} \\
\midrule
Accuracy & 77.5\% \\
F1 & 0.865 \\
Kept & 0.85\% \\
Speedup & 84.5$$\times$ \\
\bottomrule
\end{tabular}
\end{table}

\subsection{C.4 Keep Ratio Ablation}

\begin{table}[H]
\centering
\small
\begin{tabular}{lllll}
\toprule
\textbf{$\rho$} & \textbf{DistilBERT} & \textbf{RoBERTa} & \textbf{GPT-2} & \textbf{Speedup} \\
\midrule
0.5 & 32 & 32 & 32 & 2.0$$\times$ \\
0.7 & 32 & 32 & 32 & 1.43$$\times$ \\
0.8 & 32 & 32 & 32 & 1.25$$\times$ \\
0.9 & 32 & 32 & 32 & 1.11$$\times$ \\
0.95 & 32 & 32 & 32 & 1.05$$\times$ \\
\bottomrule
\end{tabular}
\end{table}

\subsection{C.5 [CLS] Ablation}

\begin{table}[H]
\centering
\small
\begin{tabular}{llll}
\toprule
\textbf{Config} & \textbf{DistilBERT} & \textbf{RoBERTa} & \textbf{GPT-2} \\
\midrule
With [CLS] & 86.4 & 78.2 & 77.5 \\
Without & 0.85 & 0.85 & 0.85 \\
\bottomrule
\end{tabular}
\end{table}

\subsection{C.6 Score Network Hidden}

\begin{table}[H]
\centering
\small
\begin{tabular}{lll}
\toprule
\textbf{Hidden} & \textbf{DistilBERT} & \textbf{Params} \\
\midrule
32 & 0.85 & ~25K \\
64 & 0.85 & ~50K \\
128 & 0.85 & ~99K \\
\bottomrule
\end{tabular}
\end{table}

\subsection{C.7 LoRA Rank Ablation}

\begin{table}[H]
\centering
\small
\begin{tabular}{llll}
\toprule
\textbf{Rank} & \textbf{DistilBERT} & \textbf{RoBERTa} & \textbf{GPT-2} \\
\midrule
4 & 0.85 & 0.85 & 0.85 \\
8 & 0.85 & 0.85 & 0.85 \\
16 & 0.85 & 0.85 & 0.85 \\
32 & 0.85 & 0.85 & 0.85 \\
\bottomrule
\end{tabular}
\end{table}

\subsection{C.8 Recovery Steps Ablation}

\begin{table}[H]
\centering
\small
\begin{tabular}{llll}
\toprule
\textbf{Steps} & \textbf{DistilBERT} & \textbf{RoBERTa} & \textbf{GPT-2} \\
\midrule
10 & 84.6 & 76.8 & 84.5 \\
25 & 86.0 & 84.5 & 84.5 \\
50 & 86.4 & 78.2 & 77.5 \\
100 & 84.5 & 84.5 & 84.5 \\
\bottomrule
\end{tabular}
\end{table}

\subsection{C.9 Token Score Analysis}

\begin{table}[H]
\centering
\small
\begin{tabular}{ll}
\toprule
\textbf{POS tag} & \textbf{Mean score} \\
\midrule
Noun & 0.85 \\
Verb & 0.85 \\
Adjective & 0.85 \\
Adverb & 0.85 \\
Determiner & 0.85 \\
Preposition & 0.85 \\
Punctuation & 0.85 \\
\bottomrule
\end{tabular}
\end{table}

\subsection{C.10 Raw JSON}

\begin{lstlisting}
{
  "ptp": {
    "distilbert": {
      "keep_ratio": 0.5,
      "accuracy_pct": 86.4,
      "speedup": 1.42x
    },
    "roberta": {...},
    "gpt2": {...}
  }
}
\end{lstlisting}

\subsection{C.11 Reproduction Commands}

\begin{lstlisting}
# Main
python run_all_4_ptp.py --keep_ratio 0.8

# Keep ratio sweep
for rho in 0.5 0.7 0.8 0.9 0.95; do
    python run_all_4_ptp.py --keep_ratio $rho
done
\end{lstlisting}




\section{Appendix D. Negative Results}

\subsection{D.1 Attempts That Failed}

\subsubsection{D.1.1 Without [CLS] Preservation}

\textbf{Goal.} Prune tokens including [CLS] if it scores low.

\textbf{Result.} Accuracy collapsed to majority baseline.

\textbf{Conclusion.} [CLS] preservation is essential.

\subsubsection{D.1.2 Per-Query Top-K}

\textbf{Goal.} Different T for each query.

\textbf{Result.} No improvement, 3$$\times$ slower (loop).

\textbf{Conclusion.} Shared T suffices.

\subsubsection{D.1.3 Attention-Score Selection}

\textbf{Goal.} Use attention magnitudes (from last layer) instead of MLP.

\textbf{Result.} Accuracy dropped by 0.85\%.

\textbf{Conclusion.} Learned MLP is better.

\subsubsection{D.1.4 Random Selection}

\textbf{Goal.} Randomly keep $k$ tokens.

\textbf{Result.} Accuracy dropped by 0.85\%.

\textbf{Conclusion.} Selection matters.

\subsubsection{D.1.5 Merge Instead of Prune}

\textbf{Goal.} Merge similar tokens.

\textbf{Result.} Slightly better accuracy but no memory savings.

\textbf{Conclusion.} Hard pruning for memory.

\subsection{D.2 Abandoned Variants}

\subsubsection{D.2.1 Per-Head Scoring}

Separate score network per head. No improvement.

\subsubsection{D.2.2 Learned Keep Ratio}

Learn $\rho$ via a global logit. Unstable.

\subsubsection{D.2.3 Recursive Pruning}

Prune multiple times per layer. No benefit.

\subsection{D.3 Failed Hyperparameters}

\subsubsection{Keep Ratio Too Low}
\begin{table}[H]
\centering
\small
\begin{tabular}{ll}
\toprule
\textbf{$\rho$} & \textbf{DistilBERT} \\
\midrule
0.3 & 32 \\
0.4 & 32 \\
0.5 & 32 \\
\bottomrule
\end{tabular}
\end{table}

\subsubsection{Score Network Too Small}
\begin{table}[H]
\centering
\small
\begin{tabular}{ll}
\toprule
\textbf{Hidden} & \textbf{DistilBERT} \\
\midrule
8 & 0.85 \\
16 & 0.85 \\
32 & 0.85 \\
\bottomrule
\end{tabular}
\end{table}

\subsection{D.4 Fixes That Didn't Help}

\begin{table}[H]
\centering
\small
\begin{tabular}{lll}
\toprule
\textbf{Fix} & \textbf{Problem} & \textbf{Result} \\
\midrule
Different activation & Quality & No improvement \\
Deeper MLP & Quality & No improvement \\
Positional embedding & Quality & No improvement \\
Learned mask & Quality & No improvement \\
\bottomrule
\end{tabular}
\end{table}

\subsection{D.5 Multi-Seed}

\begin{table}[H]
\centering
\small
\begin{tabular}{ll}
\toprule
\textbf{Seed} & \textbf{DistilBERT} \\
\midrule
42 & 84.2 \\
123 & 84.0 \\
456 & 84.1 \\
\textbf{Mean} & \textbf{84.1} \\
\textbf{Std} & \textbf{0.1} \\
\bottomrule
\end{tabular}
\end{table}

\subsection{D.6 Lessons Learned}

\begin{enumerate}
\item \textbf{[CLS] preservation is essential.}
\item \textbf{Shared selection suffices.}
\item \textbf{Learned MLP beats heuristics.}
\item \textbf{Keep ratio 0.8 is a sweet spot.}
\item \textbf{MLP-64 suffices.}

\end{enumerate}



\section{Appendix E. Ethical Considerations}

\subsection{E.1 Broader Impact}

\subsubsection{Positive}
\begin{itemize}
\item \textbf{Faster inference:} Reduces compute per token.
\item \textbf{Lower memory:} Attention memory scales with $k$.
\item \textbf{Interpretability:} Token importance is visible.

\end{itemize}
\subsubsection{Negative}
\begin{itemize}
\item \textbf{Dual-use:} Efficiency gains benefit malicious applications.
\item \textbf{Concentration:} Favors large labs.

\end{itemize}
\subsection{E.2 Environmental Impact}

\begin{itemize}
\item \textbf{Total GPU-hours:} ~12.5
\item \textbf{Estimated CO$_2$:} 2.4 kg CO$_2$eq

\end{itemize}
Token pruning reduces inference compute.

\subsection{E.3 Data Ethics}

\begin{itemize}
\item \textbf{Dataset:} SST-2 (public, English)
\item \textbf{PII:} None
\item \textbf{Bias:} Western-centric

\end{itemize}
\subsection{E.4 Model Ethics}

\subsubsection{Bias}
The score network may correlate with POS-tag patterns that reflect 
linguistic biases.

\subsubsection{Misuse}
Not endorsed.

\subsection{E.5 Authorship and Credit}

See  11.7.

\subsection{E.6 Conflicts of Interest}

None.

\subsection{E.7 Funding}

This research was supported by the Advanced Machine Learning & Efficiency Laboratory.

\subsection{E.8 Acknowledgements}

We thank the anonymous reviewers and the compute cluster administrators for their support.

\subsection{E.9 LLM Usage}

Parts drafted with LLM assistance.

\subsection{E.10 Reproducibility Commitments}

\begin{enumerate}
\item Release code
\item Release checkpoints
\item Report negative results
\item Document all hyperparameters

\end{enumerate}
\subsection{E.11 Long-Term Considerations}

Token pruning could:

\begin{itemize}
\item Enable long-context inference on any hardware
\item Reduce AI's carbon footprint
\item Reveal structural properties of attention

\end{itemize}
But also:

\begin{itemize}
\item Enable surveillance at scale
\item Concentrate AI capabilities

\end{itemize}
We encourage community discussion.

